\section{多维随机变量及其分布 II}

\subsection{条件分布}

\begin{definition}[连续型随机变量的条件分布]
    设 $(X, Y)$ 是连续型随机变量。若 $f(x, y)$ 在点 $(x, y)$ 连续，$f_Y(y)$ 在点 $y$ 连续且 $f_Y(y) > 0$，那么称
    $$
    \frac{\pdv{F}{y}}{\dv{F_Y}{y}} = \int_{- \infty}^{x} \frac{f(u, y)}{f_Y(y)} \dd{u}
    $$
    为 $Y = y$ 条件下 $X$ 的\textbf{条件分布函数}，记作 $F_{X \mid Y}(x \mid y)$。称 $\frac{f(x, y)}{f_Y(y)}$ 为 $Y = y$ 条件下 $X$ 的\textbf{条件概率密度}，记作 $f_{X \mid Y}(x \mid y)$。
\end{definition}

\subsection{随机变量的独立性}

\begin{definition}[随机变量的独立性]
    设 $(X, Y)$ 是二维随机变量，若对于任意实数 $x, y$ 都有
    $$
    \prob{X \le x \land Y \le y} = \prob{X \le x} \prob{Y \le y}
    $$
    那么称随机变量 $X, Y$ 相互独立。
\end{definition}

当 $(X, Y)$ 是连续型随机变量，并且相互独立时，其联合概率密度和边缘概率密度满足
$$
f(x, y) = f_X(x) f_Y(y)
$$
在 $\RR^2$ 上几乎处处成立。

对于离散型随机变量，独立性等价于
$$
\prob{X = x_i \land Y = y_j} = \prob{X = x_i} \prob{Y = y_j}
$$

\subsection{两个随机变量的函数的分布}

本节讨论如何由二维随机变量 $(X, Y)$ 的分布求 $Z = g(X, Y)$，以及由两个函数 $U = g_1(X, Y), V = g_2(X, Y)$ 的联合分布。

\subsubsection{离散型随机变量}

设 $(X, Y)$ 的联合分布律为
$$
p_{ij} = \prob{X = x_i \land Y = y_j}
$$
若 $Z = g(X, Y)$，对任意 $z$，有
$$
\prob{Z = z} = \sum_{(i, j): g(x_i, y_j) = z} p_{ij}
$$

\subsubsection{连续型随机变量}

设 $(X, Y)$ 的联合密度为 $f_{X,Y}(x, y)$，令 $Z = g(X, Y)$。将事件
$$
\{Z \le z\} = \{(X, Y): g(X, Y) \le z\}
$$
看作平面区域 $G_z$，则
$$
F_Z(z) = \prob{Z \le z} = \iint_{G_z} f_{X,Y}(x, y) \dd{x} \dd{y}
$$
若 $Z$ 为连续型随机变量，进一步求导得到 $f_Z(z) = F_Z'(z)$。

若 $U = g_1(X, Y), V = g_2(X, Y)$，则
$$
F_{U,V}(u, v) = \prob{U \le u \land V \le v} = \iint_{G_{u,v}} f_{X,Y}(x, y) \dd{x} \dd{y},
$$
其中 $G_{u,v} = \{(x, y): g_1(x, y) \le u,\ g_2(x, y) \le v\}$。在二阶偏导存在时，
$$
f_{U,V}(u, v) = \pdv[2]{F_{U,V}}{u}{v}
$$

\subsubsection{二维变量变换公式}

若变换 $(x, y) \mapsto (u, v)$ 存在唯一反变换
$$
x = h_1(u, v),\qquad y = h_2(u, v)
$$
且反函数连续可偏导、Jacobian 不为零，则
$$
f_{U,V}(u, v) = f_{X,Y}\ab(h_1(u, v), h_2(u, v)) \abs{\frac{\partial(x, y)}{\partial(u, v)}}
$$

\subsubsection{和的分布与卷积}

令 $Z = X + Y$。离散型情形下，
$$
\prob{Z = z} = \sum_i \prob{X = x_i \land Y = z - x_i} = \sum_j \prob{X = z - y_j \land Y = y_j}
$$
若 $X, Y$ 独立，则
$$
p_Z(z) = \sum_i p_X(x_i) p_Y(z - x_i)
$$

连续型情形下，不要求独立性即可由联合密度得到
$$
f_Z(z) = \int_{- \infty}^{+ \infty} f_{X,Y}(x, z - x) \dd{x} = \int_{- \infty}^{+ \infty} f_{X,Y}(z - y, y) \dd{y}
$$
当 $X, Y$ 独立时，$f_{X,Y}(x, y) = f_X(x) f_Y(y)$，于是
$$
f_Z(z) = \int_{- \infty}^{+ \infty} f_X(x) f_Y(z - x) \dd{x} = (f_X * f_Y)(z)
$$
这就是卷积公式。

\subsection{作业}

\begin{problem}
    习题 3.10
    \begin{solution}
        \begin{enumerate}
            \item[\textbf{1)}] $X$ 的边缘分布律为：
            \begin{center}
                \begin{mytable}{cccccc}
                    $X$ & $51$ & $52$ & $53$ & $54$ & $55$ \\
                    $\prob{X = x}$ & $0.28$ & $0.28$ & $0.22$ & $0.09$ & $0.13$ \\
                \end{mytable}
            \end{center}
            $Y$ 的边缘分布律为：
            \begin{center}
                \begin{mytable}{cccccc}
                    $Y$ & $51$ & $52$ & $53$ & $54$ & $55$ \\
                    $\prob{Y = y}$ & $0.18$ & $0.15$ & $0.35$ & $0.12$ & $0.20$ \\
                \end{mytable}
            \end{center}

            \item[\textbf{2)}] $\prob{Y = y \mid X = 51} = \frac{\prob{Y = y \land X = 51}}{\prob{X = 51}}$，于是得到：
            
            \begin{center}
                \begin{mytable}{cccccc}
                    $Y$ & $51$ & $52$ & $53$ & $54$ & $55$ \\
                    $\prob{Y = y \mid X = 51}$ & $\frac{3}{14}$ & $\frac{1}{4}$ & $\frac{5}{28}$ & $\frac{5}{28}$ & $\frac{5}{28}$ \\
                \end{mytable}
            \end{center}
        \end{enumerate}
    \end{solution}
\end{problem}

\begin{problem}
    习题 3.14
    \begin{solution}
        先求 $x, y$ 的边缘概率密度，得到：
        $$
        f_X(x) = \begin{cases}
            2 x &,\ 0 \le x \le 1 \\
            0 &,\ \text{otherwise}
        \end{cases}, \quad f_Y(y) = \begin{cases}
            1 - |y| &,\ -1 \le y \le 1 \\
            0 &,\ \text{otherwise}
        \end{cases}
        $$
        因此
        $$
        \begin{gathered}
            f_{Y \mid X}(y \mid x) = \int_{-x}^{x} \frac{f(x, u)}{f_X(x)} \dd{u} = \begin{cases}
                \frac{1}{2x} &,\ -x < y < x \\
                0 &,\ \text{otherwise}
            \end{cases}, \quad 0 < x < 1 \\
            f_{X \mid Y}(x \mid y) = \int_{|y|}^{1} \frac{f(u, y)}{f_Y(y)} \dd{u} = \begin{cases}
                \frac{1}{1 - |y|} &,\ |y| < x < 1 \\
                0 &,\ \text{otherwise}
            \end{cases}, \quad -1 < y < 1
        \end{gathered}
        $$
    \end{solution}
\end{problem}

\begin{problem}
    习题 3.18
    \begin{solution}
        \begin{enumerate}
            \item[\textbf{1)}]
            $$
            f(x, y) = \begin{cases}
                \frac{1}{2} \E^{-y / 2} &,\ 0 < x < 1 \land y > 0 \\
                0 &,\ \text{otherwise}
            \end{cases}
            $$
            \item[\textbf{2)}] 由判别式非负得到
            $$
            \Delta \ge 0 \implies X^2 \ge Y
            $$
            因此方程有实根的概率是
            $$
            \begin{aligned}
                \prob{X^2 \ge Y} & = \iint_{x^2 \ge y} f(x, y) \dd{x} \dd{y} \\
                & = \int_{0}^{1} \dd{x} \int_{x^2}^{+\infty} \frac{1}{2} \E^{-y / 2} \dd{y} \\
                & = \int_{0}^{1} \ab(1 - \E^{-x^2 / 2}) \dd{x} \\
                & = 1 - \sqrt{2\pi} \Phi(1) + \sqrt{\pi / 2}
            \end{aligned}
            $$
        \end{enumerate}
    \end{solution}
\end{problem}

\begin{problem}
    习题 3.20
    \begin{solution}
        \begin{enumerate}
            \item[\textbf{1)}] $X, Y$ 相互独立，因此
            $$
            f_{X \mid Y}(x \mid y) = f_X(x) = \begin{cases}
                \lambda \E^{-\lambda x} &,\ x > 0 \\
                0 &,\ x \le 0
            \end{cases} \quad (y > 0)
            $$

            \item[\textbf{2)}]
            $$
            \begin{aligned}
                \prob{Z = 1} & = \prob{X \le Y} = \int_{0}^{+\infty} \lambda \E^{-\lambda x} \dd{x} \int_{x}^{+ \infty} \mu \E^{-\mu y} \dd{y} \\
                & = \int_{0}^{+ \infty} \lambda \E^{-\lambda x} \cdot \E^{-\mu x} \dd{x} \\
                & = \frac{\lambda}{\lambda + \mu} \\
                \prob{Z = 0} & = \frac{\mu}{\lambda + \mu}
            \end{aligned}
            $$
            因此 $Z$ 分布律为：
            \begin{center}
                \begin{mytable}{ccc}
                    $Z$ & $0$ & $1$ \\
                    $\prob{Z = z}$ & $\frac{\mu}{\lambda + \mu}$ & $\frac{\lambda}{\lambda + \mu}$ \\
                \end{mytable}
            \end{center}
            分布函数为：
            $$
            F_Z(z) = \begin{cases}
                0 &,\ z < 0 \\
                \frac{\mu}{\lambda + \mu} &,\ 0 \le z < 1 \\
                1 &,\ z \ge 1
            \end{cases}
            $$
        \end{enumerate}
    \end{solution}
\end{problem}

\begin{problem}
    习题 3.22
    \begin{solution}
        由题可知，$f_Z = f_X * f_Y$。当 $0 \le z \le 1$ 时：
        $$
        f_Z(z) = \int_{0}^{z} 1 \cdot \E^{-(z - x)} \dd{x} = \E^{-z} \int_{0}^{z} \E^{x} \dd{x} = 1 - \E^{-z}
        $$
        当 $z > 1$ 时：
        $$
        f_Z(z) = \int_{0}^{1} 1 \cdot \E^{-(z - x)} \dd{x} = \E^{-z} \int_{0}^{1} \E^{x} \dd{x} = (\E - 1) \E^{-z}
        $$
        综上
        $$
        f_Z(z) = \begin{cases}
            1 - \E^{-z} &,\ 0 < z \le 1 \\
            (\E - 1) \E^{-z} &,\ z > 1 \\
            0 &,\ z \le 0
        \end{cases}
        $$
    \end{solution}
\end{problem}

\begin{problem}
    习题 3.24
    \begin{solution}
        \begin{enumerate}
            \item[\textbf{1)}] 分别求 $X, Y$ 的边缘概率密度，得到：
            $$
            \begin{gathered}
                f_X(x) = \begin{cases}
                    \frac{1}{2} (x + 1) \E^{-x} &,\ x > 0 \\
                    0 &,\ x \le 0
                \end{cases}, \quad f_Y(y) = \begin{cases}
                    \frac{1}{2} (y + 1) \E^{-y} &,\ y > 0 \\
                    0 &,\ y \le 0
                \end{cases}
            \end{gathered}
            $$
            因此 $f_X(x) f_Y(y) \neq f(x, y)$，得到 $X, Y$ 不独立。

            \item[\textbf{2)}] 令变换 $(Z, W) = (X + Y, X)$，那么 $(X, Y) = (W, Z - W)$，得到 Jacobian 行列式
            $$
            J = \ab|\frac{\partial(x, y)}{\partial(z, w)}| = \begin{vmatrix}
                0 & 1 \\
                1 & -1
            \end{vmatrix} = -1
            $$
            因此
            $$
            f_{Z, W}(z, w) = f(w, z - w) \cdot \abs{J} = \frac{1}{2} z \E^{-z}
            $$
            进一步得到，当 $z \ge 0$ 时：
            $$
            f_Z(z) = \int_{0}^{z} \frac{1}{2} z \E^{-z} \dd{w} = \frac{1}{2} z^2 \E^{-z}
            $$
            当 $z < 0$ 时 $f_Z(z) = 0$。
        \end{enumerate}
    \end{solution}
\end{problem}

\begin{problem}
    习题 3.25
    \begin{solution}
        由题可知 $f_Z = f_X * f_Y$。当 $z \le 2$ 时，$f_Z(z) = 0$。否则：
        $$
        f_Z(z) = \int_{1}^{z - 1} \E^{-x} \cdot \E^{-(z - x)} \dd{x} = (z - 2) \E^{-(z - 2)}
        $$
        因此
        $$
        f_Z(z) = \begin{cases}
            (z - 2) \E^{-(z - 2)} &,\ z > 2 \\
            0 &,\ z \le 2
        \end{cases}
        $$
    \end{solution}
\end{problem}

\begin{problem}
    习题 3.30
    \begin{solution}
        设寿命为 $X \sim N(160, 20^2)$，再令 $Y = \frac{X - 160}{20}$，那么 $Y \sim N(0, 1)$。根据题意，所求答案为
        $$
        P_{\mathrm{ans}} = \prob{X \ge 180}^{4} = \prob{Y \ge 1}^{4} = (1 - \Phi(1))^{4}
        $$
    \end{solution}
\end{problem}

\begin{problem}
    习题 3.34
    \begin{solution}
        $$
        \begin{aligned}
            \prob{X + Y = z} & = \sum_{k = 0}^{z} \prob{X = k} \prob{Y = z - k} \\
            & = \sum_{k = 0}^{z} \E^{-\lambda_1} \frac{\lambda_1^k}{k!} \cdot \E^{-\lambda_2} \frac{\lambda_2^{z - k}}{(z - k)!} \\
            & = \frac{\E^{-\lambda_1 - \lambda_2}}{z!} \sum_{k = 0}^{z} \binom{z}{k} \lambda_1^k \lambda_2^{z - k} \\
            & = \E^{-\lambda_1 - \lambda_2} \frac{(\lambda_1 + \lambda_2)^{z}}{z!}
        \end{aligned}
        $$
        故得证。
    \end{solution}
\end{problem}

\begin{problem}
    习题 3.35
    \begin{solution}
        $$
        \begin{aligned}
            \prob{X + Y = z} & = \sum_{k = 0}^{z} \prob{X = k} \prob{Y = z - k} \\
            & = \sum_{k = 0}^{z} \binom{n_1}{k} p^{k} (1 - p)^{n_1 - k} \cdot \binom{n_2}{z - k} p^{z - k} (1 - p)^{n_2 - (z - k)} \\
            & = p^{z} (1 - p)^{n_1 + n_2 - z} \sum_{k = 0}^{z} \binom{n_1}{k} \binom{n_2}{z - k} \\
            & = \binom{n_1 + n_2}{z} p^{z} (1 - p)^{n_1 + n_2 - z}
        \end{aligned}
        $$
        故得证。
    \end{solution}
\end{problem}